\chapter{Metodología}
\label{ch:met}

% Por redactar — estructura pendiente de definir con el asesor.

%% -----------------------------------------------------------------------
%% TODO: Contenido trasladado desde el marco teórico.
%% Descomentar y desarrollar en el siguiente entregable.
%% -----------------------------------------------------------------------

% \section{Diseño del módulo clasificador}
%
% El módulo \texttt{cv32e40p\_insn\_classifier.sv} se implementa como lógica
% combinacional pura para la clasificación y lógica secuencial para los contadores.
%
% \subsection{Integración en el ID stage}
%
% El clasificador examina \texttt{opcode} como selector primario y \texttt{funct3}
% para distinguir variantes. Por ejemplo:
% - opcode \texttt{0110011} → instrucciones registro-registro (tipo R)
%   - funct7=0000000 → aritmética (add, sll, slt, sltu, xor, srl, or, and)
%   - funct7=0100000 → aritmética (sub, sra)
% - opcode \texttt{0000011} → todos los loads (independiente de funct3)
% - opcode \texttt{1100011} → todos los branches
%
% El módulo produce seis señales de incremento de un bit, una por categoría,
% que habilitan el contador correspondiente en cada ciclo activo.
%
% \subsection{Modificación de cv32e40p\_cs\_registers.sv}
%
% La integración requiere agregar seis casos al multiplexor de lectura del
% banco de CSRs en \texttt{cv32e40p\_cs\_registers.sv}, retornando el valor
% del contador correspondiente cuando se accede a 0xBC0--0xBCB.

% -----------------------------------------------------------------------

% \section{Circuito de medición de corriente}
%
% \subsection{Topología high-side}
%
% Se adopta la topología high-side (shunt en la línea de +5V):
% - Low-side descartado: el GND del FPGA quedaría a +85mV, cualquier cable
%   externo (USB, JTAG, clip GND del osciloscopio) cortocircuitaría el shunt.
% - High-side: GND del FPGA permanece a 0V, compartido con todos los equipos.
%
% Con R_shunt = 1 Ohm y corriente típica de 85mA:
% - Caída en shunt = 85mV (1.7% de 5V) — dentro del margen ±5% de la Nexys A7
% - Jumper JP3 en posición EXT para que toda la corriente pase por el shunt
%
% \subsection{Configuración del INA240}
%
% INA240A1: G = 20 V/V
% IN+ al nodo de 5V (antes del shunt), IN- al nodo de la FPGA (después del shunt)
% → Vout = G × I × R_shunt = 20 × 85mA × 0.1Ohm = +170mV
% Alimentación: 3.3V, REF conectado a GND
%
% Resultados esperados:
% V(IN+)=5V, V(IN-)=4.9915V, V(vout)=+170mV, I=85mA

% -----------------------------------------------------------------------

% \section{Metodología de caracterización}
%
% \subsection{Bucle de instrucciones}
%
% Para cada categoría i:
% - Bucle interno (ensamblador): 1000-10000 iteraciones de una sola instrucción
% - Bucle externo (C): mantiene ejecución indefinida hasta registrar la medición
%
% \subsection{Estabilización térmica}
%
% registrar V_out en el osciloscopio. Fang [2021] observó <1% de diferencia
% entre 5 y 10 minutos de espera.
%
% \subsection{Cálculo del coeficiente}
%
% c_i = (V_out,i / (G * R_shunt)) * V_DD
%      = (V_out,i / (10 * 1)) * 5   [mW/instrucción]
%
% Se repite para las 6 categorías → vector c = [c0, c1, ..., c5]

% -----------------------------------------------------------------------

% \section{Análisis de datos en Python}
%
% line = port.readline().decode().strip()
% n = list(map(int, line.split(',')))  # n[0]..n[5]
%
% \subsection{Cálculo de potencia estimada}
%
% import numpy as np
% P_est = np.dot(c, n)   # P = sum(c_i * n_i)
%
% \subsection{Histogramas por categoría}
%
% import matplotlib.pyplot as plt
% categorias = ['Arit','Log','Mem','Branch','Jump','FP']
% plt.bar(categorias, n)
% plt.ylabel('Instrucciones ejecutadas')
% plt.title('Perfil de instrucciones')
% plt.show()
%
% \subsection{Validación del modelo}
%
% Error porcentual: epsilon = |P_est - P_real| / P_real * 100%
% P_real se mide con el osciloscopio simultáneamente a la ejecución del programa.
